\documentclass[10pt,a4paper]{amsart}
\usepackage[margin=19mm]{geometry}
\usepackage{amsmath,amssymb,mathtools}
\usepackage[scr=rsfs,cal=euler]{mathalfa}
\usepackage{xfrac}
\usepackage[unicode,hidelinks,bookmarks=false]{hyperref}
\newcommand{\ie}{i.e.\ }
\newcommand{\cf}{cf.\ }
\newcommand{\resp}{respectively\ }
\newcommand{\Subsection}{\subsection}
\providecommand{\nobreakdash}{\nobreak\hbox{-}}
\setlength{\emergencystretch}{2em}
\multlinegap0pt
\usepackage{amssymb}
\usepackage{bm}
\usepackage{float}
\usepackage{tikz}
\newtheorem{lem}{Lemma}
\newtheorem{thm}{Theorem}
\newcommand{\A}{\mathcal{A}}
\newcommand{\Ch}{\mathrm{Ch}}
\newcommand{\FF}{\mathcal{F}}
\newcommand{\Mag}{\mathrm{Mag}}
\DeclareMathOperator{\codim}{codim}
\DeclareMathOperator{\rank}{rank}
\title{Magnitude homology of real hyperplane arrangements}
\author{Junnosuke Koizumi, Ye Liu}
\date{Source: arXiv:2604.03718v4 (CC BY 4.0)}
\begin{document}
\maketitle

\noindent\emph{Source and licence.} Magnitude homology of real hyperplane arrangements. Version arXiv:2604.03718v4 (CC BY 4.0), distributed by arXiv under the Creative Commons Attribution 4.0 International licence, http://creativecommons.org/licenses/by/4.0/, as recorded in the arXiv licence metadata for this exact version and shown on the abstract page https://arxiv.org/abs/2604.03718v4. Full licence notice: \texttt{licenses/arxiv-2604.03718-CC-BY-4.0.txt}.\par\smallskip

\noindent\textit{Editorial scope.} Counted unit: the article's thm alternating (source0.tex lines 1357--1364). The article's complete original proof of it is delivered with it (source0.tex lines 1366--1441, one \texttt{proof} environment of 76 lines). No mathematics is written, repaired or re-derived: the statement and the proof are the article's own lines, in the article's own notation, and no retained line is substituted or re-flowed.  Retained uncounted as the source-local context the proof needs: lines 1341--1345.  Nothing was omitted from any retained span, so the package carries no omission record.  No retained line contains a citation, so the wrapper carries no bibliography and no bibliography entry.  No retained line contains a figure environment, an image-inclusion command or drawing code, so no image is delivered and the package declares no asset dependency.  Editorial formatting only, in addition: an AMS \texttt{amsart} preamble that restates the article's own declarations and macros used by the retained lines, namely the environments \texttt{lem}, \texttt{thm} on separate counters, together with the standard packages those lines need.  The retained environments print in this document as Lemma 1, Theorem 1 in order of appearance, whereas the article prints the same items with the numbering of its own full document.  The complete native source of the article as distributed in the arXiv e-print for this exact version is bundled unchanged as \texttt{sources/197826/source0.tex}.
\begin{lem}\label{n2positive}
    Let $\A$ be an arrangement of rank $3$.
    Let $n_k$ be the number of faces $F\in \FF(\A)$ with $\codim F=2$ and $\#\A_F=k$.
    Then we have $n_2>0$.
\end{lem}

\begin{thm}\label{alternating}
    Let $\A$ be an arrangement with $\rank \A\leq 3$.
    Let $c_\ell$ denote the coefficient of $q^\ell$ in the series expansion of $\Mag(\A)$:
    \[
    \Mag(\A)=\sum_{\ell=0}^\infty c_\ell q^\ell.
    \]
    Then we have $(-1)^\ell c_\ell\geq 0$ for all sufficiently large $\ell$.
\end{thm}

\begin{proof}
When $\rank\A\leq 2$, the claim can be checked directly.
Therefore, we may assume that $\rank \A=3$.
Let $\#\A=n$.
Write $n_k$ for the number of faces $F\in \FF(\A)$ with $\codim F=2$ and $\#\A_F=k$.
Define $(e_\ell)_{\ell\geq 0}$ by
\[
\Mag(\A,-q)=\dfrac{1}{1+(-q)^n}\biggl(\#\Ch(\A)+\sum_{k=2}^\infty 2kn_kG_k(q)\biggr)=\sum_{\ell=0}^\infty e_\ell q^\ell.
\]
It suffices to show that $e_\ell>0$ for all sufficiently large $\ell$.
Put
\[
H(q):=\#\Ch(\A)+\sum_{k=2}^{\infty}2k n_kG_k(q),\qquad
H(q)=\sum_{\ell=0}^\infty a_\ell q^\ell .
\]
We first record two elementary facts about the coefficients $g_{k,\ell}$ of $G_k$.
If $k$ is even, then
\[
G_k(q)=\frac{q(1+q^{k-1})}{(1-q)(1-q^k)}
      =\frac{q}{1-q}\sum_{j=0}^\infty\left(q^{kj}+q^{kj+k-1}\right).
\]
Hence $g_{k,\ell}$ is non-negative and non-decreasing in $\ell$.  In particular,
for $k=2$ we have $g_{2,\ell}=\ell$.
If $k$ is odd, then
\[
G_k(q)=\frac{q(1-q^{k-1})}{(1-q)(1+q^k)}
      =\frac{q(1+q+\cdots+q^{k-2})}{1+q^k}.
\]
Thus the coefficient sequence $(g_{k,\ell})_{\ell\ge 0}$ is periodic and bounded.
Since $n_2>0$ by Lemma \ref{n2positive}, there is a constant
$C>0$ such that $a_\ell\ge 4n_2\ell-C$ for all $\ell\ge 0$.

We now distinguish the parity of $n$.
First suppose that $n$ is odd.  Then $1+(-q)^n=1-q^n$, so
\[
\Mag(\A,-q)=\frac{H(q)}{1-q^n}=H(q)\sum_{j=0}^\infty q^{nj}.
\]
Therefore, with $N=\lfloor \ell/n\rfloor$,
\[
e_\ell=\sum_{j=0}^{N}a_{\ell-nj}
     \ge \sum_{j=0}^{N}\bigl(4n_2(\ell-nj)-C\bigr).
\]
The right-hand side is a polynomial of degree $1$ in $\ell$ with leading coefficient $4n_2(N+1)>0$, so it tends to $+\infty$ as $\ell\to\infty$.
Hence $e_\ell>0$ for all
sufficiently large $\ell$ in this case.

Next suppose that $n$ is even.  Then $1+(-q)^n=1+q^n$ and
\[
\frac{1}{1+q^n}=\sum_{j=0}^\infty (-1)^j q^{nj}.
\]
For even $k$, let $b_{k,\ell}$ be the coefficient of $q^\ell$ in the series expansion of $\frac{G_k(q)}{1+q^n}$.
Then we have
\[
b_{k,\ell}=\sum_{j=0}^{\lfloor \ell/n\rfloor}(-1)^jg_{k,\ell-nj}.
\]
Since $(g_{k,\ell})_\ell$ is non-negative and non-decreasing, pairing adjacent
terms gives $b_{k,\ell}\ge 0$ for every $\ell$.
For $k=2$, using $g_{2,m}=m$, we also get $b_{2,\ell}\to +\infty$.
We claim that the odd $k$ terms, after division by $1+q^n$, are
bounded.  Indeed, for odd $k$ and even $n$,
\[
\frac{G_k(q)}{1+q^n}
 =\frac{q(1+q+\cdots+q^{k-2})}{(1+q^k)(1+q^n)}.
\]
Moreover $\gcd(1+q^k,1+q^n)=1$, because $k$ is odd and $n$ is even.  Hence the
denominator has only simple roots of unity, and the coefficient sequence is
bounded.  The coefficient sequence of $\#\Ch(\A)/(1+q^n)$ is also bounded.  Thus there
is a constant $D>0$ such that the total contribution of $\#\Ch(\A)$ and all odd $k$
terms to $e_\ell$ is at least $-D$ for every $\ell$.
Combining these observations, and using Lemma \ref{n2positive}, we obtain
\[
e_\ell\ge 4n_2 b_{2,\ell}-D.
\]
Since $b_{2,\ell}\to +\infty$, the right-hand side is positive for all
sufficiently large $\ell$.  This proves the theorem.
\end{proof}

\end{document}
